% Notes-only C.50 candidate: explicit multiplier-one outer detector.
\subsection{An explicit finite outer detector from two peripheral coset modules}

\begin{statement}{C.50}{Open}{A finite characteristic outer detector with a literal recovery algorithm}
Let $\mathcal H=\ker\lambda\subset\widehat{\mathrm{GT}}_0$ and
use its canonical normalized automorphisms
\[
             \alpha_f(x)=x,\qquad \alpha_f(y)=f^{-1}yf,
                    \qquad f\in F'.                     \tag{C.50.1}
\]
Let $L$ be the explicit finite group of C.47, of order
$M=2^{293}3^{193}$ and closed derived length at most four.
There is an explicitly specified finite fully characteristic quotient
$D$ of $F/F^{(5)}$ such that
\[
       \Delta|_{\mathcal H}\quad\hbox{factors through}\quad
              \mathcal H\longrightarrow\operatorname{Out}(D).
                                                               \tag{C.50.2}
\]
The construction below gives a finite algorithm recovering the retained
C.47 conjugator residue from a representative of any outer class in
this actual image. The existing C.47 frame algorithm then evaluates
$\Delta$. No enumeration of the resulting large group is needed for
the proof, and none is claimed to have been executed.

More precisely, put $m=\exp(L)$, define
\[
 C=F\Big/\bigcap_{\psi:F\to L}\ker\psi,
                    \qquad n=|C|,
                                                               \tag{C.50.3}
\]
and take
\[
 G=\bigl(\mathbb F_2[C/\langle x_C\rangle]
          \oplus\mathbb F_2[C/\langle y_C\rangle]\bigr)\rtimes C,
 \qquad
 D=F\Big/\bigcap_{\theta:F\to G}\ker\theta.           \tag{C.50.4}
\]
All homomorphisms in these displays are continuous; every choice of
two target elements gives one. Then
\[
 n\le M^{M^2},\qquad |G|=2^{2n/m}n,
                    \qquad |D|\le |G|^{|G|^2}.        \tag{C.50.5}
\]
The exponent $m$ is the exponent of $L$, not its order.
This is an explicit multiplier-one outer detector. It neither
improves the sharp depth-four existence theorem C.48 nor computes
the smallest finite detector or the arbitrary-multiplier outer image.
\end{statement}

\begin{proof}[Complete candidate proof]
\paragraph{The actual abelianization of the first envelope.}
Write $N_L$ for the intersection in (C.50.3).
There are exactly $M^2$ homomorphisms from $F$ to $L$.
Thus $C$ is the image of their joint map into $L^{M^2}$, so it is
finite, has exponent dividing $m$, and has derived length at most
four. Precomposition with any endomorphism of $F$ permutes or
repeats maps in the intersection, proving that $N_L$ is fully
characteristic.

Let $w\in N_L$ and write its abelianization as $(a,b)\in
\widehat{\mathbb Z}^2$. For every $h\in L$, the homomorphism
$x\mapsto h,y\mapsto1$ gives $h^a=1$. Therefore
$\operatorname{ord}(h)$ divides $a$ for every $h$, and hence
$m$ divides $a$. The separate homomorphisms
$x\mapsto1,y\mapsto h$ give $m\mid b$.
Conversely every $m$th power in $F$ belongs to $N_L$, since
$L$ has exponent $m$. After taking the closed subgroup generated
by these powers and then multiplying by $F'$, these two containments
prove exactly
\[
                     C^{\rm ab}\simeq (\mathbb Z/m)^2,
       \qquad \operatorname{ord}(x_C)
              =\operatorname{ord}(y_C)=m.             \tag{C.50.6}
\]
The two generators correspond to the two standard basis vectors.
No element of $L$ itself is required to have order $m$.
This argument uses all the cyclic-subgroup homomorphisms separately,
so it also applies when the exponent is a proper least common multiple
of the orders of elements.

The quotient $C$ is exactly the universal envelope used in C.47.
That proof supplies the implication
\[
          g=(1,f)\in\mathcal H,\quad f_C=1
                         \quad\Longrightarrow\quad \Delta(g)=0.
                                                               \tag{C.50.7}
\]
We work in the original canonical marking of (C.50.1).
The change of free basis used by C.47 preserves $N_L$, since it
is fully characteristic. Thus its fixed word substitution transports
the retained residue and its evaluation algorithm to this marking.

\paragraph{A literal centralizer-coordinate equation.}
For a finite group $C$, an arbitrary element $h\in C$, and
$H=\langle h\rangle$, put $V=\mathbb F_2[C/H]$ with the
left permutation action. Let $\delta_H$ be the basis vector of
its identity coset, and set $X_h=(\delta_H,h)\in V\rtimes C$.
If $(v,c)$ centralizes $X_h$, the semidirect product law gives
\[
      ch=hc,\qquad (1-h)v=\delta_H-c\delta_H.          \tag{C.50.8}
\]
The coefficient of the left side at $H$ is zero, because $hH=H$.
The same coefficient on the right is zero if $cH=H$ and one
otherwise. These are distinct elements of the actual coefficient
field $\mathbb F_2$. Hence $c\in H$. Conversely the powers of
$X_h$ centralize it and project onto $H$. Therefore
\[
       \operatorname{im}\bigl(C_{V\rtimes C}(X_h)\to C\bigr)=H.
                                                               \tag{C.50.9}
\]
This proof imposes no condition on the order of $h$ and makes no
assertion that the affine lift has that same order. Additional
permutation-module summands do not change the displayed obstruction.

In the actual $G$ of (C.50.4), define
\[
 X=(\delta_{\langle x_C\rangle},0,x_C),\qquad
 Y=(0,\delta_{\langle y_C\rangle},y_C).
                                                               \tag{C.50.10}
\]
The free profinite universal property gives $\pi:F\to G$ with
$x\mapsto X,y\mapsto Y$. Its composite with $G\to C$ is the
canonical quotient map. Its two centralizer projections are exactly
$\langle x_C\rangle$ and $\langle y_C\rangle$ by (C.50.9).

Write $N_G$ for the intersection defining $D$. It is fully
characteristic by the same precomposition argument as above.
The homomorphism $F\to G$ obtained from the canonical map to $C$
and the zero section $C\to G$ has kernel $N_L$. Consequently
$N_G\subset N_L$, and there is an actual quotient $D\to C$.
The map $\pi$ also factors through $D$; its composite to $C$ is
this same quotient. The abelian normal subgroup in $G$ has
quotient $C$, so $G^{(5)}=1$, and hence $D^{(5)}=1$.
The two coset modules have dimensions $n/m$ by (C.50.6), proving
all the bounds in (C.50.5).

\paragraph{Outer triviality forces the retained residue to vanish.}
Suppose an actual $g=(1,f)\in\mathcal H$ acts trivially in
$\operatorname{Out}(D)$. There is $c\in D$ such that the
induced normalized automorphism equals $\operatorname{Inn}(c)$,
with $\operatorname{Inn}(c)(z)=czc^{-1}$. Its two generator
identities are
\[
             cx_Dc^{-1}=x_D,\qquad
             c y_Dc^{-1}=f_D^{-1}y_Df_D.             \tag{C.50.11}
\]
Apply the actual map $D\to G$. The first identity and the first
centralizer projection give $a=c_C\in\langle x_C\rangle$.
The second says that $f_Dc$ centralizes $y_D$, so the second
centralizer projection gives $f_Ca\in\langle y_C\rangle$.
Thus for some integers $i,j$ modulo $m$,
\[
                     a=x_C^i,\qquad f_Ca=y_C^j.       \tag{C.50.12}
\]
Since $f\in F'$, passage to (C.50.6) gives
$i(1,0)=j(0,1)$, and hence $i=j=0$ modulo $m$.
The orders in (C.50.6) now give $a=1$ and $f_C=1$ literally.
Equation (C.50.7) proves $\Delta(g)=0$.
Because $\Delta$ is a character, this kernel containment proves
its factorization through the finite actual outer image in
(C.50.2).

\paragraph{Recovery from a finite outer representative.}
For completeness, the construction gives a finite recovery procedure,
not merely an abstract kernel argument. Given
$A\in\operatorname{Aut}(D)$ representing a class in the actual
image of $\mathcal H$, enumerate $c\in D$ and $v\in D'$ until
\[
 (\operatorname{Inn}(c)\circ A)(x_D)=x_D,\qquad
 (\operatorname{Inn}(c)\circ A)(y_D)=v^{-1}y_Dv.       \tag{C.50.13}
\]
At least one pair exists: the actual canonical representative and its
actual conjugator residue supply it. The image $v_C$ is independent
of every choice in this finite search. To prove this, let two choices
give normalized representatives $\alpha_v,\alpha_w$ with
$\alpha_w=\operatorname{Inn}(d)\circ\alpha_v$.
The identities on $x_D,y_D$ say that $d$ centralizes $x_D$
and that $wdv^{-1}$ centralizes $y_D$. Applying the two literal
centralizer projections gives
\[
                  d_C=x_C^i,\qquad w_Cd_Cv_C^{-1}=y_C^j.
                                                               \tag{C.50.14}
\]
Both $v,w$ belong to $D'$, so their images in $C^{\rm ab}$
vanish. Equation (C.50.6) again forces $i=j=0$, whence
$w_C=v_C$. Comparing with the actual canonical representative
proves that this is its actual retained conjugator residue.

Apply the fixed C.47 marking substitution to this residue, project
$C$ to the specific affine quotient $L$ used there, and run its
existing finite coset-module frame solver. It returns the exact
value of $\Delta$. The procedure is asserted on the actual outer
image only; it is not an algorithm recognizing which arbitrary
finite outer classes lift to the full gentle group. Neither this
construction nor the finite recovery proof asserts an
arbitrary-multiplier normalization or evaluation.
\end{proof}
